Die bisherige Beschreibung des Elektrons\index{Elektron} beruhte weitgehend auf
klassischen Vorstellungen. Das Elektron besitzt eine Masse\index{Masse}, eine
elektrische Ladung\index{Ladung!elektrische} und bewegt sich in elektrischen
und magnetischen Feldern\index{elektrisches Feld}\index{magnetisches Feld}
nach bestimmten Kraftgesetzen. Mit diesen Vorstellungen lassen sich viele
Experimente erklären, zum Beispiel die Ablenkung von Elektronenstrahlen
\index{Elektronenstrahl} oder ihre Bewegung in magnetischen Feldern.

Zu Beginn des 20. Jahrhunderts zeigte sich jedoch, dass diese klassische
Beschreibung nicht ausreicht. Elektronen verhalten sich nicht immer wie kleine
geladene Teilchen\index{Teilchen}, die sich auf eindeutig bestimmbaren Bahnen
\index{Bahn} bewegen. Unter bestimmten Bedingungen zeigen sie Eigenschaften,
die man sonst von Wellen\index{Welle} kennt. Besonders deutlich wird dies bei
Beugungs-\index{Beugung} und Interferenzexperimenten\index{Interferenz}.

Damit wird das Elektron zu einem zentralen Objekt der Quantenmechanik
\index{Quantenmechanik}. Sein Zustand\index{Zustand!quantenmechanischer}
wird nicht mehr durch eine klassische Bahn beschrieben, sondern durch eine
Wellenfunktion\index{Wellenfunktion}. Diese Wellenfunktion liefert keine direkte
Bahnkurve, sondern Wahrscheinlichkeiten\index{Wahrscheinlichkeit} für mögliche
Messergebnisse\index{Messung}. Das Elektron bleibt ein Teilchen, aber seine
Beschreibung erhält einen wellenartigen Charakter.

Dieser Gedanke führt zur de-Broglie-Hypothese\index{de-Broglie-Hypothese},
nach der jedem Teilchen mit Impuls\index{Impuls} eine Wellenlänge
\index{Wellenlänge} zugeordnet werden kann. Beim Elektron ist dieser
Wellencharakter\index{Wellencharakter} experimentell nachweisbar. Der
Doppelspaltversuch\index{Doppelspaltversuch} mit Elektronen zeigt besonders
eindrucksvoll, dass auch einzelne Elektronen ein Interferenzmuster
\index{Interferenzmuster} aufbauen können, wenn viele Messungen zusammen
betrachtet werden.

Die mathematische Grundlage dieser Beschreibung bildet die
Schrödingergleichung\index{Schrödingergleichung}. Für das Wasserstoffatom
\index{Wasserstoffatom} liefert sie nicht mehr feste Elektronenbahnen
\index{Elektronenbahn} wie im bohrschen Atommodell
\index{bohrsches Atommodell}, sondern diskrete Zustände
\index{Zustand!diskreter} und räumliche Wahrscheinlichkeitsverteilungen
\index{Wahrscheinlichkeitsverteilung}. Aus diesen Lösungen entsteht die
Vorstellung der Atomorbitale\index{Atomorbital}\index{Orbital}.

Zusätzlich besitzt das Elektron eine weitere rein quantenmechanische
Eigenschaft: den Spin\index{Spin}. Zusammen mit dem Pauli-Prinzip
\index{Pauli-Prinzip} erklärt er, warum Elektronen in Atomen bestimmte
Zustände besetzen und warum die Elektronenhülle\index{Elektronenhülle}
eine geordnete Struktur besitzt.

In diesem Kapitel wird das Elektron daher nicht mehr als klassisches Teilchen
behandelt, sondern als quantenmechanisches Objekt\index{quantenmechanisches Objekt}.
Diese Sichtweise ist die Grundlage für das moderne Verständnis von Atomen
\index{Atom}, chemischen Bindungen\index{chemische Bindung}, Festkörpern
\index{Festkörper} und vielen elektronischen Anwendungen
\index{Anwendung!elektronische}.

